\chapter{DF-053 Audit motion and shader effects family across real surfaces}\label{record-098}

\noindent{\footnotesize\color{BookMuted}\path{docs/design/DF-053-audit-motion-and-shader-effects-family-across-real-surfaces.md}\par}\smallskip

\begin{RecordMeta}
title: DF-053 Audit motion and shader effects family across real surfaces
legend: DF
lane: design
priority: medium
keywords:
  - component-audit
  - dogfood
  - lowerings
  - motion

\end{RecordMeta}

\section{Framing}\label{record-098-framing}

Motion and shader effects have become a real authoring surface, not a demo side quest. \texttt{c\allowbreak{}a\allowbreak{}n\allowbreak{}v\allowbreak{}a\allowbreak{}s\allowbreak{}(\allowbreak{})} now has quad, Braille, and glyph-fit resolution modes, the raytrace helpers are available for app-owned shader scenes, and spring animation is stable under slow terminal pulses.

The DOGFOOD docs still treat that family as a small animated field. That is too thin for current repo truth. The story should prove the full family against real runtime behavior instead of trusting isolated fixtures or prose alone.

\section{Sponsored Users}\label{record-098-sponsored-users}

\begin{itemize}
\tightlist
\item
  TUI app authors who want expressive terminal visuals without inventing their own canvas, ray, or animation substrate.
\item
  Documentation readers evaluating whether motion, shader, and transition effects lower honestly across rich, static, pipe, and accessible output.
\item
  Maintainers who need regression tests that keep examples, docs, and exported APIs aligned as the rendering stack evolves.
\end{itemize}

\section{Hills}\label{record-098-hills}

\begin{enumerate}
\def\labelenumi{\arabic{enumi}.}
\tightlist
\item
  A TUI app author can inspect one DOGFOOD story and see \texttt{c\allowbreak{}a\allowbreak{}n\allowbreak{}v\allowbreak{}a\allowbreak{}s\allowbreak{}(\allowbreak{})} cover quad, Braille, glyph-fit, and raytrace-backed shader output without local hacks.
\item
  A docs reader can switch the same story across rich, static, pipe, and accessible profiles and get truthful lowerings instead of decorative glyphs leaking into constrained modes.
\item
  A maintainer can run a focused cycle test that proves \texttt{c\allowbreak{}a\allowbreak{}n\allowbreak{}v\allowbreak{}a\allowbreak{}s\allowbreak{}(\allowbreak{})}, transition shaders, \texttt{a\allowbreak{}n\allowbreak{}i\allowbreak{}m\allowbreak{}a\allowbreak{}t\allowbreak{}e\allowbreak{}(\allowbreak{})}/spring motion, and \texttt{t\allowbreak{}i\allowbreak{}m\allowbreak{}e\allowbreak{}l\allowbreak{}i\allowbreak{}n\allowbreak{}e\allowbreak{}(\allowbreak{})} remain represented by real examples and DOGFOOD previews.
\end{enumerate}

\section{Playback Questions}\label{record-098-playback-questions}

\begin{itemize}
\tightlist
\item
  Does the motion/shader DOGFOOD story include variants for shader waves, Braille fields, glyph-fit raytraced surfaces, and spring/timeline motion orchestration?
\item
  Do pipe and accessible previews describe the state-change meaning without relying on Braille, block, or ornamental shader characters?
\item
  Does a static profile render a deterministic visual snapshot that still exercises the shader path?
\item
  Do example sources for transitions and spring animation compile against the current public API surface?
\end{itemize}

\section{Requirements}\label{record-098-requirements}

\begin{itemize}
\tightlist
\item
  Keep the existing motion/shader story identity and coverage family id.
\item
  Add real story variants instead of only expanding prose.
\item
  Use the current \texttt{c\allowbreak{}a\allowbreak{}n\allowbreak{}v\allowbreak{}a\allowbreak{}s\allowbreak{}(\allowbreak{})} API, including \texttt{r\allowbreak{}e\allowbreak{}s\allowbreak{}o\allowbreak{}l\allowbreak{}u\allowbreak{}t\allowbreak{}i\allowbreak{}o\allowbreak{}n\allowbreak{}:\allowbreak{} \allowbreak{}'\allowbreak{}g\allowbreak{}l\allowbreak{}y\allowbreak{}p\allowbreak{}h\allowbreak{}'}, for at least one DOGFOOD preview.
\item
  Use exported raytrace helpers inside the DOGFOOD preview so the raytrace kernel is proven from a real docs surface.
\item
  Represent spring motion and \texttt{t\allowbreak{}i\allowbreak{}m\allowbreak{}e\allowbreak{}l\allowbreak{}i\allowbreak{}n\allowbreak{}e\allowbreak{}(\allowbreak{})} together in a deterministic preview.
\item
  Keep pipe and accessible lowerings textual and meaning-first.
\item
  Repair any stale example code uncovered by the audit.
\end{itemize}

\section{Acceptance Criteria}\label{record-098-acceptance-criteria}

\begin{itemize}
\tightlist
\item
  \texttt{t\allowbreak{}e\allowbreak{}s\allowbreak{}t\allowbreak{}s\allowbreak{}/\allowbreak{}c\allowbreak{}y\allowbreak{}c\allowbreak{}l\allowbreak{}e\allowbreak{}s\allowbreak{}/\allowbreak{}D\allowbreak{}F\allowbreak{}-\allowbreak{}0\allowbreak{}5\allowbreak{}3\allowbreak{}/\allowbreak{}m\allowbreak{}o\allowbreak{}t\allowbreak{}i\allowbreak{}o\allowbreak{}n\allowbreak{}-\allowbreak{}s\allowbreak{}h\allowbreak{}a\allowbreak{}d\allowbreak{}e\allowbreak{}r\allowbreak{}-\allowbreak{}d\allowbreak{}o\allowbreak{}g\allowbreak{}f\allowbreak{}o\allowbreak{}o\allowbreak{}d\allowbreak{}-\allowbreak{}a\allowbreak{}u\allowbreak{}d\allowbreak{}i\allowbreak{}t\allowbreak{}.\allowbreak{}t\allowbreak{}e\allowbreak{}s\allowbreak{}t\allowbreak{}.\allowbreak{}t\allowbreak{}s} fails before the DOGFOOD story and stale example are updated.
\item
  The motion/shader story exposes \texttt{s\allowbreak{}h\allowbreak{}a\allowbreak{}d\allowbreak{}e\allowbreak{}r\allowbreak{}-\allowbreak{}w\allowbreak{}a\allowbreak{}v\allowbreak{}e}, \texttt{b\allowbreak{}r\allowbreak{}a\allowbreak{}i\allowbreak{}l\allowbreak{}l\allowbreak{}e\allowbreak{}-\allowbreak{}f\allowbreak{}i\allowbreak{}e\allowbreak{}l\allowbreak{}d}, \texttt{g\allowbreak{}l\allowbreak{}y\allowbreak{}p\allowbreak{}h\allowbreak{}-\allowbreak{}r\allowbreak{}a\allowbreak{}y\allowbreak{}t\allowbreak{}r\allowbreak{}a\allowbreak{}c\allowbreak{}e}, and \texttt{s\allowbreak{}p\allowbreak{}r\allowbreak{}i\allowbreak{}n\allowbreak{}g\allowbreak{}-\allowbreak{}t\allowbreak{}i\allowbreak{}m\allowbreak{}e\allowbreak{}l\allowbreak{}i\allowbreak{}n\allowbreak{}e} variants.
\item
  Rendering every motion/shader variant in \texttt{i\allowbreak{}n\allowbreak{}t\allowbreak{}e\allowbreak{}r\allowbreak{}a\allowbreak{}c\allowbreak{}t\allowbreak{}i\allowbreak{}v\allowbreak{}e}, \texttt{s\allowbreak{}t\allowbreak{}a\allowbreak{}t\allowbreak{}i\allowbreak{}c}, \texttt{p\allowbreak{}i\allowbreak{}p\allowbreak{}e}, and \texttt{a\allowbreak{}c\allowbreak{}c\allowbreak{}e\allowbreak{}s\allowbreak{}s\allowbreak{}i\allowbreak{}b\allowbreak{}l\allowbreak{}e} story profiles produces non-empty output.
\item
  Pipe and accessible variants avoid Braille and block-art glyph leakage.
\item
  The spring example no longer passes removed \texttt{a\allowbreak{}n\allowbreak{}i\allowbreak{}m\allowbreak{}a\allowbreak{}t\allowbreak{}e\allowbreak{}(\allowbreak{})} options.
\end{itemize}

\section{Implementation Outline}\label{record-098-implementation-outline}

\begin{itemize}
\tightlist
\item
  Expand \texttt{m\allowbreak{}o\allowbreak{}t\allowbreak{}i\allowbreak{}o\allowbreak{}n\allowbreak{}P\allowbreak{}r\allowbreak{}e\allowbreak{}v\allowbreak{}i\allowbreak{}e\allowbreak{}w\allowbreak{}(\allowbreak{})} so it can render wave, Braille, glyph-raytrace, and spring/timeline previews from one helper.
\item
  Import the needed public \texttt{b\allowbreak{}i\allowbreak{}j\allowbreak{}o\allowbreak{}u\allowbreak{}-\allowbreak{}t\allowbreak{}u\allowbreak{}i} helpers through the package surface rather than private module paths.
\item
  Add DF-053 cycle tests that render the DOGFOOD story through the story protocol instead of snapshotting internals.
\item
  Remove stale spring-example drift discovered during the audit.
\end{itemize}

\section{Drift Check}\label{record-098-drift-check}

\begin{itemize}
\tightlist
\item
  \texttt{e\allowbreak{}x\allowbreak{}a\allowbreak{}m\allowbreak{}p\allowbreak{}l\allowbreak{}e\allowbreak{}s\allowbreak{}/\allowbreak{}d\allowbreak{}o\allowbreak{}c\allowbreak{}s\allowbreak{}/\allowbreak{}s\allowbreak{}t\allowbreak{}o\allowbreak{}r\allowbreak{}i\allowbreak{}e\allowbreak{}s\allowbreak{}.\allowbreak{}t\allowbreak{}s} now imports \texttt{b\allowbreak{}i\allowbreak{}j\allowbreak{}o\allowbreak{}u\allowbreak{}-\allowbreak{}t\allowbreak{}u\allowbreak{}i} from the workspace source barrel so DOGFOOD tests read current package truth instead of stale ignored \texttt{d\allowbreak{}i\allowbreak{}s\allowbreak{}t\allowbreak{}/} artifacts.
\item
  \texttt{e\allowbreak{}x\allowbreak{}a\allowbreak{}m\allowbreak{}p\allowbreak{}l\allowbreak{}e\allowbreak{}s\allowbreak{}/\allowbreak{}s\allowbreak{}p\allowbreak{}r\allowbreak{}i\allowbreak{}n\allowbreak{}g\allowbreak{}/\allowbreak{}m\allowbreak{}a\allowbreak{}i\allowbreak{}n\allowbreak{}.\allowbreak{}t\allowbreak{}s} still demonstrated \texttt{a\allowbreak{}n\allowbreak{}i\allowbreak{}m\allowbreak{}a\allowbreak{}t\allowbreak{}e\allowbreak{}(\allowbreak{})} with a removed \texttt{f\allowbreak{}p\allowbreak{}s} option; the audit repaired that example while keeping the same runtime story.
\item
  No follow-on backlog item is needed for this slice. The remaining broader opportunity is to add a dedicated visual capture for the motion/shader story, but the current tests already prove static, pipe, accessible, and rich-story rendering.
\end{itemize}

\section{Playback}\label{record-098-playback}

\begin{itemize}
\tightlist
\item
  RED: \texttt{t\allowbreak{}e\allowbreak{}s\allowbreak{}t\allowbreak{}s\allowbreak{}/\allowbreak{}c\allowbreak{}y\allowbreak{}c\allowbreak{}l\allowbreak{}e\allowbreak{}s\allowbreak{}/\allowbreak{}D\allowbreak{}F\allowbreak{}-\allowbreak{}0\allowbreak{}5\allowbreak{}3\allowbreak{}/\allowbreak{}m\allowbreak{}o\allowbreak{}t\allowbreak{}i\allowbreak{}o\allowbreak{}n\allowbreak{}-\allowbreak{}s\allowbreak{}h\allowbreak{}a\allowbreak{}d\allowbreak{}e\allowbreak{}r\allowbreak{}-\allowbreak{}d\allowbreak{}o\allowbreak{}g\allowbreak{}f\allowbreak{}o\allowbreak{}o\allowbreak{}d\allowbreak{}-\allowbreak{}a\allowbreak{}u\allowbreak{}d\allowbreak{}i\allowbreak{}t\allowbreak{}.\allowbreak{}t\allowbreak{}e\allowbreak{}s\allowbreak{}t\allowbreak{}.\allowbreak{}t\allowbreak{}s} failed on the missing story variants, incomplete related-family docs, and stale spring example option.
\item
  GREEN: the motion/shader story now renders \texttt{s\allowbreak{}h\allowbreak{}a\allowbreak{}d\allowbreak{}e\allowbreak{}r\allowbreak{}-\allowbreak{}w\allowbreak{}a\allowbreak{}v\allowbreak{}e}, \texttt{b\allowbreak{}r\allowbreak{}a\allowbreak{}i\allowbreak{}l\allowbreak{}l\allowbreak{}e\allowbreak{}-\allowbreak{}f\allowbreak{}i\allowbreak{}e\allowbreak{}l\allowbreak{}d}, \texttt{g\allowbreak{}l\allowbreak{}y\allowbreak{}p\allowbreak{}h\allowbreak{}-\allowbreak{}r\allowbreak{}a\allowbreak{}y\allowbreak{}t\allowbreak{}r\allowbreak{}a\allowbreak{}c\allowbreak{}e}, and \texttt{s\allowbreak{}p\allowbreak{}r\allowbreak{}i\allowbreak{}n\allowbreak{}g\allowbreak{}-\allowbreak{}t\allowbreak{}i\allowbreak{}m\allowbreak{}e\allowbreak{}l\allowbreak{}i\allowbreak{}n\allowbreak{}e} variants across every documented profile.
\item
  The glyph raytrace preview uses \texttt{c\allowbreak{}a\allowbreak{}n\allowbreak{}v\allowbreak{}a\allowbreak{}s\allowbreak{}(\allowbreak{})} with \texttt{r\allowbreak{}e\allowbreak{}s\allowbreak{}o\allowbreak{}l\allowbreak{}u\allowbreak{}t\allowbreak{}i\allowbreak{}o\allowbreak{}n\allowbreak{}:\allowbreak{} \allowbreak{}'\allowbreak{}g\allowbreak{}l\allowbreak{}y\allowbreak{}p\allowbreak{}h\allowbreak{}'} and the exported raytrace helpers.
\item
  The spring/timeline preview uses the exported motion timeline and spring stepping helpers to produce a deterministic frame.
\item
  Pipe and accessible previews lower to explicit text without Braille or block glyph leakage.
\end{itemize}

\section{Retrospective}\label{record-098-retrospective}

\begin{itemize}
\tightlist
\item
  The real drift was not in \texttt{c\allowbreak{}a\allowbreak{}n\allowbreak{}v\allowbreak{}a\allowbreak{}s\allowbreak{}(\allowbreak{})} this time; it was DOGFOOD breadth. The docs surface had not caught up with the rendering primitives already landed.
\item
  The source-barrel import is important for this repo because \texttt{d\allowbreak{}i\allowbreak{}s\allowbreak{}t\allowbreak{}/} is ignored and can lag behind source during active cycles.
\item
  Testing the story protocol directly is a better guard than static prose checks: it proves variants render through the same path the docs app uses.
\end{itemize}

